\chapter{Phase 1: Baseline Reimplementation}
\section{Purpose}
Phase 1 separates implementation verification from reproduction of the full AACOPF method. The first subphase uses a known approximate initial pose and a conventional bootstrap particle filter. Only after this baseline behaves correctly should unknown-pose initialization and the AACOPF-specific particle transition mechanism be introduced.

\section{Experiment P1.1--P1.4: deterministic DR, drift, UWB and PF}
A deterministic curved trajectory is generated at \SI{10}{Hz}. Noise-free dead reckoning must reproduce the ground truth to numerical precision. Then distance and heading-increment perturbations are introduced. A moving auxiliary trajectory supplies noisy range observations. The PF uses the same corrupted DR increments and the UWB likelihood.

\section{Current results}
The numerical values in this section are generated by \texttt{experiments/run\_phase1.py} from \texttt{configs/phase1.yaml}. They are a verification run, not a statistical performance claim. For seed 42, ideal DR reproduces the ground truth exactly within floating-point precision. Noisy DR reaches a position RMSE of \SI{1.058}{m} and a final position error of \SI{2.820}{m}. PF+UWB reduces the position RMSE to \SI{0.338}{m} and the final error to \SI{0.682}{m}. The corresponding heading RMSE decreases from \SI{7.00}{deg} to \SI{2.54}{deg}.

\begin{figure}[ht]
\centering
\begin{tikzpicture}
\begin{axis}[width=.82\textwidth,height=.55\textwidth,xlabel={$x$ [m]},ylabel={$y$ [m]},axis equal image,grid=both,legend pos=north west]
\addplot table[x=x_true,y=y_true,col sep=comma]{../results/phase1/trajectory_report.csv};\addlegendentry{Ground truth}
\addplot table[x=x_dr,y=y_dr,col sep=comma]{../results/phase1/trajectory_report.csv};\addlegendentry{Noisy DR}
\addplot table[x=x_pf,y=y_pf,col sep=comma]{../results/phase1/trajectory_report.csv};\addlegendentry{PF+UWB}
\end{axis}
\end{tikzpicture}
\caption{Ground-truth trajectory, noisy dead reckoning, and PF+UWB estimate.}
\end{figure}

\section{Observability diagnostic}
For the state $[p_x,p_y,\psi]^\top$, the local transition and range Jacobians are
\begin{equation}
F_k=\begin{bmatrix}1&0&-\Delta L_k\sin\psi_k\\0&1&\Delta L_k\cos\psi_k\\0&0&1\end{bmatrix},\quad
H_k=\begin{bmatrix}(p_x-a_x)/d&(p_y-a_y)/d&0\end{bmatrix}.
\end{equation}
A finite-horizon local observability matrix is assembled and its singular values recorded. With the stationary auxiliary trajectory used in the diagnostic, the singular values are approximately $(176.17,10.07,4.94\times10^{-15})$, indicating a rank-deficient local linearization. With the selected moving auxiliary trajectory they are approximately $(110.39,11.93,6.15)$, giving a smallest-to-largest singular-value ratio of $5.57\times10^{-2}$. This is consistent with the geometric symmetry argument and with the role of auxiliary motion discussed by Han et al.~\cite{han2020cooperative}. It is a finite-horizon local diagnostic, not by itself a global nonlinear observability proof.

\section{Interpretation}
This first run establishes the intended software invariant: ideal DR is exact, noisy DR accumulates drift, and a conventional PF using moving-node UWB ranges substantially reduces that drift. The result is deliberately modest: it verifies the fusion pipeline before attempting a more difficult paper reproduction. The next phase-1 steps are (i) multi-seed verification, (ii) unknown initial pose with annular position/heading hypotheses, (iii) explicit stationary-versus-moving auxiliary convergence tests, and only then (iv) an audited AACOPF implementation following Han et al.~\cite{han2020cooperative}.
